% Compile with XeLaTeX or Tectonic. Scientific data are generated by prepare_data.py.
\documentclass[UTF8,a4paper,zihao=-4,fontset=none]{ctexart}
\usepackage[margin=23mm,top=24mm,bottom=25mm,headheight=16pt]{geometry}
\usepackage{fontspec}
\setCJKmainfont{Noto Serif CJK SC}
\setCJKsansfont{Noto Sans CJK SC}
\setCJKmonofont{Noto Sans Mono CJK SC}
\usepackage{amsmath,amssymb,booktabs,tabularx,longtable,array}
\usepackage{graphicx,xcolor,tikz,pgfplots}
\usetikzlibrary{arrows.meta,positioning,calc}
\usepgfplotslibrary{fillbetween}
\pgfplotsset{compat=1.18,table/col sep=comma}
\usepackage{enumitem,fancyhdr,lastpage,caption,placeins}
\usepackage[numbers,sort&compress]{natbib}
\usepackage{xurl}
\usepackage[unicode,colorlinks=true,linkcolor=blue!45!black,citecolor=blue!45!black,urlcolor=blue!45!black]{hyperref}
\hypersetup{pdftitle={连接组约束的果蝇全脑模型中习惯化样神经响应的最小慢机制},pdfauthor={作者待确认；Digital Fly Brain 项目研究报告},pdfsubject={固定连接与局部短时突触抑制的比较；计算研究报告，未经同行评审},pdfkeywords={果蝇,连接组,习惯化,短时突触抑制,可重复计算}}
\definecolor{Ink}{HTML}{213547}
\definecolor{Blue}{HTML}{2166AC}
\definecolor{Orange}{HTML}{C45A20}
\definecolor{Gray}{HTML}{555555}
\linespread{1.22}
\setlength{\parskip}{0.15em}
\setlength{\emergencystretch}{2em}
\setcounter{tocdepth}{1}
\hypersetup{bookmarksdepth=2}
\renewcommand{\topfraction}{0.9}
\renewcommand{\textfraction}{0.08}
\renewcommand{\floatpagefraction}{0.8}
\setlist{itemsep=2pt,topsep=4pt,leftmargin=2em}
\captionsetup{font=small,labelfont=bf,labelsep=quad}
\renewcommand{\arraystretch}{1.16}
\ctexset{section/format=\Large\bfseries,subsection/format=\large\bfseries}
\pagestyle{fancy}
\fancyhf{}
\fancyhead[L]{\small 连接组约束的习惯化样神经响应}
\fancyhead[R]{\small 研究报告\enspace v1.0}
\fancyfoot[C]{\small \thepage\,/\,\pageref*{LastPage}}
\renewcommand{\headrulewidth}{0.3pt}
\newenvironment{primer}[1]{\begin{quote}\small\noindent\textbf{\color{Blue}读者导读：#1}\par\smallskip}{\end{quote}}
\newcolumntype{Y}{>{\raggedright\arraybackslash}X}
\newcommand{\ms}{\,\mathrm{ms}}
\newcommand{\secunit}{\,\mathrm{s}}
\newcommand{\mV}{\,\mathrm{mV}}
\newcommand{\Mzero}{\mathrm{M0}}
\newcommand{\Mone}{\mathrm{M1}}
\newcommand{\ReportAuthor}{作者与单位待确认}
% Primitive input avoids LaTeX file-hook tokens becoming a spurious table row.
\makeatletter
\newcommand{\inputrows}[1]{\@@input #1 }
\makeatother
\title{\vspace{-1.0em}\bfseries 连接组约束的果蝇全脑模型中\\习惯化样神经响应的最小慢机制\\[0.5em]
\large 固定连接与局部短时突触抑制的可重复计算比较\\[0.6em]
\normalsize\normalfont A Minimal Slow Mechanism for Habituation-like Neural Responses\\in a Connectome-Constrained Drosophila Whole-Brain Model}
\author{\ReportAuthor\\\small Digital Fly Brain 项目研究报告}
\date{\small 实验日期：2026 年 9 月 16 日\quad 报告版本：1.0}

\begin{document}
\maketitle
\thispagestyle{fancy}
\begin{center}\small\color{Gray}
论文体例的研究草稿；尚未投稿或同行评审。结论限于所测试的计算模型与协议。
\end{center}
\begin{abstract}
\noindent\textbf{背景：}真实连接图谱能够约束感觉信息如何传播，但它是否足以产生跨次刺激的反应递减，不能仅凭网络规模判断。\textbf{目的：}检验原始固定连接果蝇全脑模型是否表现习惯化样神经响应，并评估一个局部慢状态机制的条件性充分性。\textbf{方法：}在 Shiu 等人的 FlyWire v630 模型上保留 127,400 个神经元和 14,687,178 条聚合连接，向 70 个 JO-CE 机械感觉神经元提供重复的 200 ms 随机驱动。比较固定连接模型 M0 与仅在 4,764 条兴奋性输出边上增加短时突触抑制的 M1；采用两种刺激起始间隔、五种冻结输入实现和训练末态的独立恢复分支。\textbf{结果：}20 条序列共 280 个响应窗口。M0 的 aBN1 计数逐次恒定；M1 在 0.5 s 间隔下的平均早晚期递减为 90.22\%，休息 10 s 后组均值恢复至首次水平；2 s 间隔下平均递减为 9.88\%。实际感觉脉冲序列在同一种子内完全一致。M1 同时降低首次响应，部分下游读出存在地板效应。\textbf{结论：}在所测条件下，局部可恢复资源状态足以产生递减、恢复与部分间隔依赖，而原始固定连接动力学没有表现该效应。研究尚未确立真实动物机制、完整习惯化判据或连接组相对于简单模型的独特优势。

\medskip\noindent\textbf{关键词：}果蝇；连接组；习惯化；泄漏积分发放模型；短时突触抑制；可重复计算
\end{abstract}

\clearpage
\section*{English abstract}
\addcontentsline{toc}{section}{English abstract}
\begin{small}
\noindent\textbf{Background.} Anatomical connectivity constrains signal propagation but does not by itself establish memory across repeated stimuli.
\textbf{Objective.} We tested whether a fixed-connectivity whole-brain model exhibits habituation-like neural responses and whether local short-term synaptic depression is sufficient to generate selected features.
\textbf{Methods.} We retained 127,400 neurons and 14,687,178 aggregated directed edges of the Shiu et al. FlyWire v630 model. Frozen 200-ms stochastic input patterns were replayed to 70 JO-CE mechanosensory neurons. The static model (M0) was compared with a model containing one recovery timescale on 4,764 excitatory afferent output edges (M1; depletion fraction 0.01, recovery time constant 2 s). Five input realizations, two onset intervals, and independent 2-s and 10-s recovery branches yielded 20 sequences and 280 response windows.
\textbf{Results.} Repeated aBN1 spike counts were constant in M0. In M1, the mean early-to-late decrement was 90.22\% at a 0.5-s interval and 9.88\% at a 2-s interval. At 10-s rest, group-mean responses returned to the initial M1 level, although individual realizations did not recover identically. Sensory-neuron spike trains were identical within each seed across conditions. Reduced initial descending-neuron responses exposed a floor-effect limitation.
\textbf{Conclusions.} Local recoverable resource dynamics are sufficient for selected habituation-like neural features under the tested conditions. These results neither identify the mechanism in living flies nor establish stimulus specificity, dishabituation, or a unique computational benefit of the full connectome. This is a reproducible computational pilot study, not a validated behavioral model.
\end{small}

\section*{面向非专业读者的摘要}
\addcontentsline{toc}{section}{面向非专业读者的摘要}
持续听到无害的背景声时，我们常会越来越少地注意它，但休息一段时间后又可能重新作出反应。这类“对重复而无关紧要的刺激减少反应”的现象称为习惯化。本报告没有模拟人的注意力，也没有直接观察果蝇的动作，而是问了一个更小、可检验的问题：\textbf{把同一串神经信号反复送入数字果蝇的大脑模型，某个相关神经元会不会越来越少地发放信号？}

答案取决于模型中是否保存了足够持久的状态。原模型在本实验中每次反应相同。我们在一小组输入连接上加入“使用后暂时减少、休息后补充”的资源变量后，快速重复刺激能够引起明显递减，休息能够使其恢复。这个结果好比证明了“带有可充电缓冲器的线路可以做出这种现象”，\textbf{并没有证明真实果蝇一定使用这个缓冲器，也没有证明整个大脑才能做到它}。

建议第一次阅读时先看第\ref{sec:background}节、图\ref{fig:design}和第\ref{sec:results}节；数学与软件细节可随后阅读。正文中的“读者导读”提供直觉解释，附录给出术语、原始逐种子结果与复现路径。

\vspace{0.4em}
\begin{center}\small\color{Gray}
证据层级：文献事实\quad$\vert$\quad 本项目实测\quad$\vert$\quad 机制解释\quad$\vert$\quad 后续建议。四者在文中分别表述。
\end{center}
\clearpage
\tableofcontents
\clearpage
\section{引言：连接图谱并不等于会学习的大脑}
\label{sec:intro}
成年果蝇脑的突触级连接组为计算建模提供了少见的解剖约束：我们不必从随机连接开始，而可以使用真实神经元之间的连接方向及接触数量\citep{dorkenwald2024}。Shiu 等人据此构建全脑泄漏积分发放模型，展示了味觉相关活动、感觉通路之间的相互作用以及触角梳理相关回路的可预测性\citep{shiu2024}。然而，单次刺激能够沿正确通路产生响应，不代表该模型已经拥有跨刺激记忆。

习惯化适合检验这个差别。与复杂任务训练相比，它不要求语言、显式奖赏或长期参数优化，却要求系统的当前响应依赖最近的刺激历史\citep{rankin2009}。如果一个模型每次收到同样输入都作出同样响应，它可以完成感觉传递，但没有在相应协议下表现出反应递减。相反，若引入某个慢变量后出现下降，也不能仅凭一条曲线确定动物使用了相同机制。

本研究提出两个递进问题：\textbf{（Q1）}原始固定连接的全脑动力学，在重复 JO-CE 刺激和足够长的间隔下，是否产生跨次响应递减与恢复？\textbf{（Q2）}如果没有，加入局部、可恢复的短时突触资源状态，能否在保留真实连接结构的条件下产生这些特征？本轮并未直接检验更强的\textbf{（Q3）}：真实连接组是否比简单动力学模型或随机结构具有不可替代的预测优势。

本报告的贡献是建立连续状态、输入严格配对且可审计的计算基线，记录一个明确的阴性结果和一个有局限的充分性结果。\textbf{短时突触抑制不是本研究新提出的生物学机制；“最小”指本轮只加入一种候选慢机制，不是已经证明全局数学最小性。}

\section{背景与概念：给没有神经科学基础的读者}
\label{sec:background}
\subsection{神经元、突触与连接组分别是什么}
神经元可以用膜电位表示其电活动状态。真实神经元的机制很复杂，本模型把它简化为一个会积累输入、同时向静息状态泄漏的单元。膜电位超过阈值时，单元发出一次离散事件，称为脉冲或动作电位；随后状态被重置，并短暂进入不应期。

突触是一个神经元影响另一个神经元的连接。\emph{兴奋性}影响在本模型中增加后续发放的倾向；\emph{抑制性}影响则减少这种倾向。它们不是“好”和“坏”的连接：抑制有助于选择响应、稳定回路，并可能参与习惯化。两个神经元之间可以有多个解剖突触，本模型把这些接触聚合为一条带权有向边。因此，\textbf{神经元数、聚合边数和解剖突触数是三个不同的数量}。

连接组回答“哪些细胞连接到哪些细胞”，但通常不直接提供所有膜时间常数、突触释放动力学或动物当时的内部状态。把连接组转化为一个可运行模型，还必须作出关于这些动力学的假设。解剖连接也不应直接等同于被测动物完整的生前记忆。

\begin{primer}{地图、交通规则与实时车流}
可以把连接组类比为道路地图，神经元和突触方程类比为交通规则，随时间改变的膜电位、脉冲和资源量则类似实时车流。有了地图，不代表交通规则已经确定；规则固定，也不代表车流没有历史依赖。这个类比只帮助理解结构与状态的区别，不是大脑与道路的机制等同。
\end{primer}

\subsection{为什么“不改长期权重”也可能保留记忆}
本研究区分三类量。第一类是\emph{结构}，如边是否存在、突触数量和连接方向；第二类是\emph{参数}，如阈值、时间常数和总体权重尺度；第三类是\emph{动态状态}，如当前膜电位、剩余突触驱动或可用资源。即使结构和参数保持不变，状态仍可随过去的刺激而变化，并影响下一次输出。

与典型机器学习中的优化不同，本轮没有梯度下降、训练损失或逐边权重拟合。M1 中的有效传递强度确实会随资源改变，因此它包含短时可塑性；但这不等于做了长期权重训练。本文所说的“12 次训练刺激”是习惯化实验的暴露序列，不是运行一个机器学习优化器。

原始 LIF 网络本来就有动态状态和循环连接。不能从“没有专门的可塑性规则”推出“固定权重网络绝不可能习惯化”。循环网络也可能形成慢活动模式。我们报告的是实际测试范围内的阴性结果，而不是对所有固定权重网络的否定。

\subsection{习惯化、感觉适应与疲劳为什么不能混为一谈}
习惯化通常指重复刺激导致的反应减弱，且不能简单归因于感觉器官越来越不灵敏或运动器官失去执行能力\citep{rankin2009}。观察到输出下降至少有四种可能：输入真的变弱了；外围感觉器官适应了；中枢处理保留了历史；输出通路失能了。因此，“越来越少发放”只是现象，不是机制诊断。

经典研究列出了反应递减、休息恢复、频率依赖、刺激特异性和去习惯化等特征。所有系统不必同时满足全部特征，但研究者应清楚说明检验了哪些。尤其要区分两种测试：训练 A 后对 B 仍有响应，属于特异性或泛化问题；插入 B 后对原刺激 A 的响应恢复，才是去习惯化测试。后者还需要排除休息本身带来的恢复以及普遍敏化。

本轮把相同感觉脉冲直接回放给模型，控制了刺激编码的变化；同时检测了神经读出的休息恢复。它没有真实感觉器官和完整身体，因此\textbf{不能声称已在生物实验意义上排除全部感觉适应或运动疲劳解释}。

\subsection{相关工作如何约束模型选择}
最小动力学研究表明，一个会积累并衰减的内部状态，加上非线性读出，就可能生成多项习惯化特征；因此，全脑规模并不是产生下降和恢复曲线的必要前提\citep{smart2024}。这种认识提高了连接组研究的验收标准：真正有价值的下一步是预测刺激变化和神经干预，而不只是解释一条已经看过的曲线。

生物实验还提供了不同的候选机制。果蝇嗅觉习惯化研究支持局部 GABA 能抑制性传递的可塑性；在相应实验中，短期习惯化涉及约 30 分钟的气味暴露及分钟级恢复\citep{das2011}。味觉研究支持抑制性神经元相关机制，并进一步研究了新刺激如何解除已经形成的习惯化\citep{paranjpe2012,trisal2022}。这些结果支持把慢抑制和去抑制列为竞争模型，但不能直接证明 JO-CE 也依赖相同机制。

视觉 light-off 惊跳和其他果蝇范式具有丰富的习惯化研究基础\citep{engel2009}；近期个体动力学建模进一步说明，低维状态模型具有相当的解释能力\citep{boon2026}。本轮没有把任意视觉激活当作 light-off，也没有把秒级神经输出替代分钟级嗅觉或味觉行为。

\subsection{为什么选 JO-CE，而不是继续局限于糖刺激}
Johnston 器官是果蝇触角中的机械感觉器官。JO-CE 是其中一组感觉神经元；aBN1 是触角梳理相关中间神经元，aDN1 和 aDN2 是相关下行神经元。原全脑研究已经给出这些细胞的可用标识，且 JO-CE 激活 aBN1 的预测得到生理验证\citep{shiu2024}。这为正常反应提供了依据，允许先讨论“重复之后发生什么”。

选择该通路是本轮可操作性与科学可证伪性的折中，并非认定它比嗅觉或视觉更重要。输入被定义为神经脉冲，不等价于已校准的声音、振动波形或触角触碰。JO-F 在原模型中不能可靠激活 aBN1，因此不能简单把 JO-F 当作一个可直接比较的“新刺激 B”。最重要的是，aBN1 的活动只是神经读出，不是果蝇梳理动作本身。

\section{材料与方法}
\label{sec:methods}
\subsection{研究设计与证据来源}
研究方案和正式协议在正式结果产生前保存于本地。其形式属于预先规定的本地方案，不是公开注册平台上的预注册。工程预试验使用一个种子、三个重复刺激和一个恢复分支；通过后没有依据曲线调整正式输入频率、脉冲宽度、STD 参数、主读出或判据。

正式实验比较两种模型、两个刺激起始间隔和五个输入实现。每条序列含 12 次重复刺激，随后从同一训练末态派生两个独立恢复分支。所有条件都使用同一个连接组。研究流程见图\ref{fig:design}；完整协议为 \path{habituation/protocols/main_v1.json}。

\begin{figure}[!htb]
\centering
\small\textbf{A. 模型与读出：解剖结构保留，候选慢机制局部改变传递}\par\medskip
\resizebox{0.97\linewidth}{!}{%
\begin{tikzpicture}[>=Latex,font=\footnotesize,
 box/.style={draw=Ink!65,rounded corners=2pt,align=center,minimum height=1.0cm,minimum width=2.4cm,fill=Ink!3}]
\node[box] (input) at (0,0) {冻结随机输入\\逐次相同回放};
\node[box] (jon) at (3.1,0) {70 个 JO-CE\\机械感觉细胞};
\node[box,minimum width=3.1cm] (brain) at (6.7,0) {保留完整连接网络\\127,400 个神经元};
\node[box,minimum width=3.1cm] (readout) at (10.8,0) {aBN1（主要读出）\\aDN1/aDN2（次要读出）};
\draw[->,thick] (input)--(jon);
\draw[->,thick] (jon)--(brain);
\draw[->,thick] (brain)--(readout);
\node[align=center,text=Blue] (std) at (4.8,-1.5)
{M0：固定传递强度\\M1：JO-CE 兴奋性输出边\\引入可恢复资源状态};
\draw[->,Blue] (std.north)--(4.5,-0.1);
\node[box,dashed,fill=white,minimum height=0.85cm] (body) at (10.8,-1.65) {完整身体与梳理动作\\本轮未建模};
\draw[->,dashed] (readout)--(body);
\end{tikzpicture}}

\medskip\textbf{B. 连续训练与独立恢复分支（示意，不按时间比例绘制）}\par\medskip
\resizebox{0.95\linewidth}{!}{%
\begin{tikzpicture}[>=Latex,font=\footnotesize,
 pulse/.style={draw=Blue,fill=Blue!9,rounded corners=2pt,align=center,minimum width=1.1cm,minimum height=0.8cm},
 state/.style={draw=Ink,align=center,minimum width=1.8cm,minimum height=0.9cm}]
\node[pulse] (p1) at (0,0) {刺激 1};
\node[pulse] (p2) at (1.7,0) {刺激 2};
\node (ellipsis) at (3.1,0) {$\cdots$};
\node[pulse] (p12) at (4.5,0) {刺激 12};
\node[state,fill=Ink!5] (snap) at (7,0) {同一训练末态\\保存快照};
\draw[->] (p1)--node[above] {$T$}(p2);
\draw[->] (p2)--(ellipsis);\draw[->] (ellipsis)--(p12);\draw[->] (p12)--(snap);
\node[state,minimum width=3.0cm] (r2) at (11,0.8) {休息 2 s\\$\rightarrow$ 探测 A};
\node[state,minimum width=3.0cm] (r10) at (11,-0.8) {休息 10 s\\$\rightarrow$ 探测 A};
\draw[->] (snap.east)--++(0.5,0)|-(r2.west);
\draw[->] (snap.east)--++(0.5,0)|-(r10.west);
\node[align=center] at (2.6,-1.05) {刺激宽度 200 ms；响应窗口 300 ms\\$T=0.5$ s 或 2 s；序列内部不重置};
\end{tikzpicture}}
\caption{实验逻辑示意。A 中箭头表示输入和测量流程，不是完整解剖回路图；JO-CE 与被监测细胞均属于同一全脑网络，图中方框不是互斥细胞集合。M1 改变 JO-CE 的全部选中兴奋性输出边，不限于直接连到 aBN1 的边。B 中两个探测分支独立恢复同一训练末态，因此短休息探测不会污染长休息结果。}
\label{fig:design}
\end{figure}


\subsection{连接图谱与模型边界}
固定上游实现为 Shiu 等人的公开代码\citep{shiu2024}，提交为
\begin{center}\small\texttt{91bdd1e7dcf193f3e7ca5a8933497fcef63b7960}\end{center}
采用 FlyWire v630 处理后的细胞和连接表：127,400 个神经元、14,687,178 条聚合有向边，对应 52,793,639 个解剖突触。没有为降低成本而裁剪子图；仿真期间保留所有细胞和静态连接。这里的“全脑”沿用该模型范围，不包括虚拟身体、完整腹神经索或全部生理过程。

70 个 JO-CE 输入 ID 从固定版本 Notebook 的字面量清单提取，不执行整个 Notebook。主要读出为 aBN1，次要描述性读出为 aDN1/aDN2；精确标识列于附录表\ref{tab:ids}。本文并未对原论文的全部实验进行完整复现。

\subsection{M0：原始固定连接 LIF 模型}
对非不应期中的神经元 $j$，状态方程为
\begin{align}
\tau_m\frac{\mathrm{d}v_j}{\mathrm{d}t}&=v_0-v_j+g_j,\label{eq:lif}\\
\tau_g\frac{\mathrm{d}g_j}{\mathrm{d}t}&=-g_j.\label{eq:g}
\end{align}
$v_j$ 为膜电位，$g_j$ 为上游实现中的有效突触驱动，\textbf{单位同样是电压，不是电导}。当 $v_j>v_{\mathrm{th}}$ 时发出脉冲，将 $v_j$ 置为 $v_{\mathrm{rst}}$、$g_j$ 置零，并进入不应期。上游实现把上述两个微分方程都标记为在不应期内停止更新；本研究没有更改该约定。数值积分使用 Brian2 的线性状态更新方法及 Cython 后端\citep{stimberg2019}。

来自细胞 $i$ 的脉冲经过固定传递延迟后，令
\begin{equation}
g_j\leftarrow g_j+w_{ij},\qquad
w_{ij}=n_{ij}s_iw_{\mathrm{syn}},\label{eq:weight}
\end{equation}
其中 $n_{ij}$ 为聚合突触数量，$s_i$ 是上游传递符号约定，$w_{\mathrm{syn}}$ 为全局尺度。正式实验保持上游静态参数不变，见表\ref{tab:params}。固定连接的含义是没有显式、随刺激历史变化的突触效能规则；它不表示 $v$ 和 $g$ 没有历史。

\begin{table}[tbp]
\centering\small
\caption{固定模型参数与本研究新增参数。M1 的两个 STD 参数为预先选定的演示设置，并非 JO-CE 生理拟合结果。}
\label{tab:params}
\begin{tabularx}{\linewidth}{lllY}
\toprule
类别 & 符号 & 数值 & 含义\\
\midrule
静态 LIF & $v_0,v_{\mathrm{rst}}$ & $-52\mV$ & 静息与重置电位\\
& $v_{\mathrm{th}}$ & $-45\mV$ & 发放阈值\\
& $\tau_m$ & $20\ms$ & 膜时间常数\\
& $\tau_g$ & $5\ms$ & 突触驱动衰减时间\\
& $t_{\mathrm{rfc}}$ & $2.2\ms$ & 普通细胞不应期；被驱动输入细胞设为零\\
& $t_{\mathrm{dly}}$ & $1.8\ms$ & 突触传递延迟\\
& $w_{\mathrm{syn}}$ & $0.275\mV$ & 每个解剖突触的全局权重尺度\\
数值/输入 & $\Delta t$ & $0.1\ms$ & 积分及输入事件网格\\
& $f$ & $150\,\mathrm{Hz}$ & 刺激期间的等效单细胞输入率\\
新增 STD & $U$ & $0.01$ & 每次到达事件的资源消耗比例\\
& $\tau_{\mathrm{rec}}$ & $2\secunit$ & 资源恢复时间常数\\
\bottomrule
\end{tabularx}
\end{table}

\begin{primer}{怎样理解时间常数}
时间常数不是“多久后绝对归零”，而是指数变化的速度尺度。一个单纯衰减的量经过一个时间常数后约剩原来的 37\%。原模型的显式单细胞时间常数为毫秒级，而新增资源状态的恢复尺度为秒级。不过，不能仅凭这个差别排除网络产生慢模式的可能，所以仍需要实际运行 M0。
\end{primer}

\subsection{M1：JO-CE 兴奋性输出的局部短时突触抑制}
\label{sec:std}
M1 只替换以 JO-CE 为起点、静态权重为正的 4,764 条聚合输出边。对每条选中连接引入无量纲可用资源 $x_{ij}\in[0,1]$，初始为 1。事件之间满足
\begin{equation}
\frac{\mathrm{d}x_{ij}}{\mathrm{d}t}=\frac{1-x_{ij}}{\tau_{\mathrm{rec}}}.
\label{eq:resource}
\end{equation}
当一个突触前脉冲到达时，先用到达前资源传递，再消耗资源：
\begin{align}
g_j&\leftarrow g_j+w_{ij}x_{ij}^{-},\\
x_{ij}^{+}&=(1-U)x_{ij}^{-}.\label{eq:deplete}
\end{align}
这里的“抑制”指短时传递效能下降（depression），\textbf{不是把这些兴奋性连接变成 GABA 能抑制性连接}。该机制也不同于“加强一组抑制性神经元”的竞争假说。

恢复采用事件驱动的精确指数更新。对于相邻到达事件，先把上一次事件后的状态推进到当前时刻，再执行式\eqref{eq:deplete}。仿真实现中将对应原始静态边的权重置零，另建同起点、同终点、同延迟的动态边，避免双重传递；其他连接不变。它没有新增解剖通路，也没有按刺激序号人为乘一个越来越小的包络。

\begin{primer}{暂时用掉的“资源”是什么}
可以把 $x$ 想成一个使用后需要补充的缓冲器：$x=1$ 时传递完整，$x=0.5$ 时同一个脉冲在这条边上只产生一半有效驱动。真实生物中可能有多个过程导致类似现象；本研究的 $x$ 是抽象状态，不是直接测到的囊泡数，也没有验证其一定对应 JO-CE 的真实生理过程。
\end{primer}

“一个恢复时间尺度”不等于“全脑只有一个状态变量”。当前实现按边存储资源，所有边共享 $U$ 和 $\tau_{\mathrm{rec}}$。在相同初态、相同延迟与参数下，同一突触前细胞对应的这些边具有相同资源更新历史，理论上可以共享状态；本轮没有实施该优化。两项共享参数没有通过数据拟合，4,764 条边也不是 4,764 个被独立训练的参数。

\subsection{输入：冻结随机实现与严格配对}
每个输入神经元在刺激持续的 200 ms 内，在每个 $0.1\ms$ 时间格独立抽取一次 Bernoulli 事件，概率为
\begin{equation}
p=f\Delta t=150\times 10^{-4}=0.015.
\end{equation}
这是小步长的离散稀疏随机驱动，不是严格连续时间 Poisson 过程。单次刺激的理论输入事件均值为 $70\times150\times0.2=2,100$。随机种子为 1701--1705。

每个种子的事件阵列只生成一次，然后在该种子的所有重复刺激、两个模型、两个间隔和恢复探测中回放。因此，同一序列内不会因“下一次碰巧更少的输入”而产生假递减。跨种子改变输入实现，以检查结果对有限输入变化的稳健性。

事件通过强数值驱动注入输入神经元的膜电位，幅度为上游的 $250w_{\mathrm{syn}}=68.75\mV$；与上游一样，将这些被驱动细胞的不应期设为零。这个注入是促使模型细胞发放的数值接口，\textbf{不是对真实机械刺激的生理幅度校准}。注入事件仍需经过 LIF 阈值与重置调度，不能未经检查就假定每个事件必然一对一变成感觉脉冲。分析器因此比较实际感觉神经元发放的逐事件序列，而不只比较注入计数。

\subsection{连续状态刺激与独立恢复分支}
\label{sec:protocol}
第 $k$ 次刺激（$k=1,\ldots,12$）的起始时刻为 $(k-1)T$，$T\in\{0.5,2\}\secunit$。刺激持续 0.2 s；响应在起始后的 0.3 s 窗口内计数。因此，两个协议的\emph{刺激结束到下次刺激开始}静默间隔分别为 0.3 s 和 1.8 s，不能把 $T$ 本身误写为静默间隔。

\textbf{整个重复序列内不重建网络，不清空膜电位、突触驱动、不应期或待传递事件。}仅在切换独立模型条件、输入实现或刺激间隔时恢复初态。完整训练段结束时刻是 $11T+0.3\secunit$，短/长间隔分别为 5.8 s 和 22.3 s。

训练末态保存后，分支 A 等待 2 s 再给一次原刺激，分支 B 重新恢复\emph{同一个末态}、等待 10 s 后再给原刺激。休息从末次 300 ms 响应窗口结束开始计，而不是从末次刺激结束计；所以相对于末次刺激结束，两个探测间隔实际为 2.1 s 和 10.1 s。这样的分支设计避免了 2 s 探测本身改变 10 s 结果。

\subsection{读出、效应量与预先规定的筛查}
\label{sec:metrics}
对每个种子、模型和刺激间隔，令 $R_k$ 为第 $k$ 次刺激后 300 ms 内 aBN1 的脉冲数。定义
\begin{equation}
I=R_1,\qquad E=\frac{R_1+R_2+R_3}{3},\qquad
L=\frac{R_{10}+R_{11}+R_{12}}{3},\qquad
D=1-\frac{L}{E}.\label{eq:D}
\end{equation}
其中 $I$ 是首次响应，$E$ 为早期均值，$L$ 为末期均值，$D$ 为递减比例。$E=0$ 时不定义该比值，不能把无初始响应当作完全习惯化。本文先逐种子计算 $D$，再平均；它通常不等于先汇总 $L,E$ 后再作比值。

H1 的实用性筛查阈值为 $D\geq0.30$。在 H1 通过的前提下，H2 筛查要求 10 s 探测 $P_{10}$ 同时满足
\begin{equation}
P_{10}-L\geq0.20E,\qquad P_{10}\geq0.80I.\label{eq:H2}
\end{equation}
平台筛查要求第 7--9 次与第 10--12 次响应均值之差不超过 $0.10E$。这些是本项目预先规定的描述性门槛，不是文献对习惯化的普遍数值定义，也不是统计显著性检验。

对恢复另报告 $P_t/I$。看过部分结果后增加的探索性诊断是 $(P_t-L)/(I-L)$，仅在 $I>L$ 时定义；它衡量“已经损失的响应恢复了多少”，与“回到首次水平的多少”不同。该探索性指标没有被事后加入 H4 的通过判据。

\subsection{统计单位、工程对照与软件证据}
五个种子是同一连接组上的五种输入实现，\textbf{不是五只动物}。一条序列内的 12 个响应也不是 12 个独立生物样本。结果以逐种子值、均值、极值范围和筛查通过数描述；没有构造生物总体置信区间或宣称显著性。

正式实验前，M1-zero 将 $U$ 设为零，检查替换突触的实现是否与 M0 等价；三个 pilot 模型分别通过训练末态恢复重放检查。正式实验用独立分析入口回读原始事件，校验文件哈希、细胞 ID 精度、时间网格、重复事件、资源边界、刺激时序和窗口计数，再验证 20 条序列的完整覆盖。

使用 Python 3.10.21、Brian2 2.5.1、NumPy 1.22.3、Pandas 1.4.3、Cython 0.29.33 与 PyArrow 17.0.0；主仿真单进程运行，限制原生数学库线程。完整环境沿用既有独立复现环境，未修改系统 Python 或上游源码。软件、参数和数据的固定方式见附录\ref{app:repro}。
\FloatBarrier

\section{结果}
\label{sec:results}
\subsection{工程门槛与正式数据覆盖}
预试验中，M0 的三个 aBN1 响应依次为 10、10、10，2 s 休息后为 10；M1-zero 相同，且两者的完整全脑训练脉冲表一致。M1 的对应响应为 8、5、3，休息后为 6。三个模型各自保存并恢复训练末态后，恢复分支的完整脉冲表均能一致重放。预试验表明实现能够保留状态，并通过“零强度扩展不改变结果”的工程门槛；它没有被当作独立生物验证集。

正式实验完整覆盖 20 条序列，含 240 个训练响应和 40 个恢复响应，共 280 个窗口。重新回读 60 个原始事件文件后，各读出计数、总脉冲数及感觉脉冲数均与登记结果一致。同一种子的感觉神经元相对发放时刻，在模型、间隔、重复刺激和恢复探测之间完全相同。登记文件中的窗口外脉冲数为零；这不是程序把状态清零的结果，连续状态仍由方程推进。

\subsection{M0 在两个间隔下均未产生逐次递减}
图\ref{fig:training}与表\ref{tab:main}显示，固定连接模型的 aBN1 响应逐次恒定。种子 1701--1704 的每次计数均为 10，种子 1705 均为 9；两种间隔下组均值均为 9.80，早期和末期相等，所有 $D=0$。2 s、10 s 探测同样保持各自的初始计数。

aDN1 与 aDN2 也在各序列内保持恒定。这个结果不是“不够显著的下降”，而是在当前确定性回放输入和所报告计数精度下，根本没有观察到跨次计数变化。它支持一个有限结论：\textbf{在所测试的 JO-CE 输入、原始参数和两个间隔下，M0 没有表现出本研究定义的响应递减。}

\newcommand{\TrainingPanel}[2]{%
\begin{tikzpicture}
\begin{axis}[width=\linewidth,height=5.6cm,title={#2},
 xlabel={刺激序号},ylabel={aBN1 脉冲数},xmin=0.7,xmax=12.3,ymin=-0.3,ymax=11,
 xtick={1,3,6,9,12},ytick={0,2,4,6,8,10},grid=major,grid style={black!10},
 tick label style={font=\footnotesize},label style={font=\small},title style={font=\small\bfseries}]
\addplot[name path=staticlow,draw=none,forget plot] table[x=trial,y=low]{data/train_static_#1.csv};
\addplot[name path=statichigh,draw=none,forget plot] table[x=trial,y=high]{data/train_static_#1.csv};
\addplot[Gray!15,forget plot] fill between[of=staticlow and statichigh];
\addplot[name path=stdlow,draw=none,forget plot] table[x=trial,y=low]{data/train_std_#1.csv};
\addplot[name path=stdhigh,draw=none,forget plot] table[x=trial,y=high]{data/train_std_#1.csv};
\addplot[Blue!15,forget plot] fill between[of=stdlow and stdhigh];
\addplot[Gray,thick,mark=square*,mark size=1.5pt] table[x=trial,y=mean]{data/train_static_#1.csv};
\addplot[Blue,thick,mark=*,mark size=1.6pt] table[x=trial,y=mean]{data/train_std_#1.csv};
\end{axis}
\end{tikzpicture}}
\begin{figure}[tbp]
\centering
\begin{minipage}{0.49\linewidth}\TrainingPanel{500}{A. 起始间隔 0.5 s}\end{minipage}\hfill
\begin{minipage}{0.49\linewidth}\TrainingPanel{2000}{B. 起始间隔 2 s}\end{minipage}
\par\smallskip{\small\textcolor{Gray}{$\blacksquare$ M0：固定连接}\qquad\textcolor{Blue}{$\bullet$ M1：局部 STD}}
\caption{重复刺激下 aBN1 的响应。每点为五个冻结输入实现的均值；阴影为最小值至最大值，不是置信区间。计数窗口为每次刺激起始后的 300 ms。两幅图纵轴相同。M0 逐次恒定；M1 在短间隔下强烈递减，在长间隔下保留较大响应。横轴是刺激序号，不是相同长度的物理时间。}
\label{fig:training}
\end{figure}


\begin{table}[tbp]
\centering\small
\setlength{\tabcolsep}{5pt}
\caption{主要读出 aBN1 的结果。除 $D$ 外单位均为每个 300 ms 窗口的脉冲数；每格为五个输入实现的均值。$I,E,L$ 分别为首次、前 3 次和后 3 次响应，$P_t$ 为休息 $t$ 后的独立探测。$D$ 先按每个种子计算再平均。}
\label{tab:main}
\begin{tabular}{lrrrrrr}
\toprule
模型 / 间隔 & $I$ & $E$ & $L$ & $D$（\%） & $P_2$ & $P_{10}$\\
\midrule
\inputrows{data/main_rows.tex}
\bottomrule
\end{tabular}
\end{table}

\subsection{M1 在短间隔下产生强递减与休息恢复}
在 $T=0.5\secunit$ 条件下，M1 首次响应均值为 8.20，前 3 次均值为 5.40，后 3 次均值降至 0.53。逐种子递减比例均值为 90.22\%，范围为 81.25\%--100\%；5/5 输入实现通过 H1 筛查，也全部通过平台筛查。

短间隔组 2 s 探测均值为 5.40，10 s 探测均值为 8.20（图\ref{fig:recovery}）。所有五例都满足预设 H2 门槛。但组均值恢复不意味着每个实现逐点完全恢复：10 s 探测相对首次的比值为 87.5\%--112.5\%。这些略低或略高于首次的计数不能被删去或截断为“100\%”；本文保留其原值，不用单一资源指数直接解释所有网络输出细节。

还应注意，M1 首次响应已经低于 M0 的 9.80，均值降幅约为 16.3\%。因此，该机制不仅储存跨刺激历史，也影响第一段 200 ms 刺激内部的传递。若只展示每组各自归一化后的下降曲线，就可能掩盖这项正常功能损失。

\begin{primer}{90\% 下降到底相对什么}
这里的 90.22\% 是“后 3 次平均响应，相对前 3 次平均响应”的下降；不是相对未加 STD 的 M0，也不是相对第一次。第一例短间隔的 $E=(8+5+3)/3=5.33$、$L=1$，所以 $D=1-1/5.33=81.25\%$。先算每例再平均，得到组结果。这个明确的分母决定了百分数的意义。
\end{primer}

\newcommand{\RecoveryPanel}[2]{%
\begin{tikzpicture}
\begin{axis}[width=\linewidth,height=5.7cm,title={#2},ylabel={aBN1 脉冲数},xmin=-0.4,xmax=3.4,ymin=-0.3,ymax=11,
 xtick={0,1,2,3},xticklabels={首次,末期,休息2秒,休息10秒},ytick={0,2,4,6,8,10},
 grid=major,grid style={black!10},tick label style={font=\scriptsize},label style={font=\small},title style={font=\small\bfseries}]
\addplot[Gray,only marks,mark=square*,mark size=2pt,error bars/.cd,y dir=both,y explicit]
 table[x expr=\thisrow{stage}-0.07,y=mean,y error minus=minus,y error plus=plus]{data/recovery_static_#1.csv};
\addplot[Blue,only marks,mark=*,mark size=2pt,error bars/.cd,y dir=both,y explicit]
 table[x expr=\thisrow{stage}+0.07,y=mean,y error minus=minus,y error plus=plus]{data/recovery_std_#1.csv};
\end{axis}
\end{tikzpicture}}
\begin{figure}[tbp]
\centering
\begin{minipage}{0.49\linewidth}\RecoveryPanel{500}{A. 训练间隔 0.5 s}\end{minipage}\hfill
\begin{minipage}{0.49\linewidth}\RecoveryPanel{2000}{B. 训练间隔 2 s}\end{minipage}
\par\smallskip{\small\textcolor{Gray}{$\blacksquare$ M0：固定连接}\qquad\textcolor{Blue}{$\bullet$ M1：局部 STD}}
\caption{首次、末期及独立恢复探测的响应。点为五个输入实现的均值，误差棒为 min--max；“末期”先在每个实现内平均第 10--12 次，再跨实现汇总。横轴是四个离散条件，\textbf{不是均匀时间轴}，因此不连接为恢复曲线。10 s 时 M1 组均值回到首次水平，但个别实现低于或高于首次。两个恢复点不足以估计完整恢复动力学。}
\label{fig:recovery}
\end{figure}


\subsection{刺激间隔改变效应强度，但尚不能确定恢复速率规律}
在 $T=2\secunit$ 条件下，M1 的首次均值同样为 8.20，早期均值为 7.33，末期均值为 6.60；平均 $D$ 只有 9.88\%，范围为 0--18.18\%。五例均未达到 30\% 的 H1 门槛，4/5 通过平台筛查。较频繁刺激引起更强递减的方向在五个配对输入实现中一致（图\ref{fig:paired}），支持本范围内的 H4(a) 描述。

长间隔组的 2 s、10 s 探测均值分别为 6.40、8.20。不能因为 10 s 时两组都接近初始均值，就声称恢复时间常数相同；同样不能因为短间隔恢复的绝对增量更大，就声称其恢复更快。两组训练末态的损失深度不同，长间隔还有一例未通过平台标准，而且仅有两个恢复点。

经典 H4(b) 关心达到渐近水平后的频率依赖恢复。当前数据没有建立这一特征，也不足以严格否定它。M1 中所有局部资源确实共享同一个规定的恢复时间常数，但神经输出还要经过循环网络和阈值非线性；资源恢复速率不能不经检验地替代神经输出恢复速率。

\begin{figure}[tbp]
\centering
\begin{tikzpicture}
\begin{axis}[width=10.2cm,height=5.8cm,xmin=0.7,xmax=2.3,ymin=-4,ymax=106,
 xtick={1,2},xticklabels={0.5 s 间隔,2 s 间隔},ytick={0,20,40,60,80,100},ylabel={递减比例 $D$（\%）},
 grid=major,grid style={black!10},tick label style={font=\small},label style={font=\small}]
\pgfplotstableread{data/paired_decrement.csv}\pairedtable
\pgfplotsinvokeforeach{0,1,2,3,4}{%
 \pgfplotstablegetelem{#1}{D500}\of\pairedtable\edef\firstvalue{\pgfplotsretval}
 \pgfplotstablegetelem{#1}{D2000}\of\pairedtable\edef\secondvalue{\pgfplotsretval}
 \addplot[Blue!70,thick,mark=*,mark size=2pt] coordinates {(1,\firstvalue) (2,\secondvalue)};
}
\addplot[Gray,dashed,thick,mark=square*,mark size=2pt] coordinates {(1,0) (2,0)};
\addplot[Orange,densely dotted,forget plot] coordinates {(0.75,30) (2.25,30)};
\node[anchor=west,font=\footnotesize,text=Orange] at (axis cs:0.75,34) {预设 30\% 筛查阈值};
\end{axis}
\end{tikzpicture}
\caption{同一种子内的配对间隔比较。五条蓝线分别对应 M1 的五个输入实现；灰色虚线表示所有 M0 实现的 $D=0$。每条线连接同一输入在两种刺激间隔下的描述性效应量，不表示一个连续的参数拟合曲线。短间隔的递减在五个实现中都大于长间隔。部分端点重叠；具体数值见附录表\ref{tab:individual}。}
\label{fig:paired}
\end{figure}


\subsection{下游读出的地板效应限制了行为解释}
表\ref{tab:downstream}显示，M1 不仅改变 aBN1，也降低了相关下行输出。短间隔末期的 aDN1 和 aDN2 平均计数都降至零。但 aDN1 的 M1 首次平均已经只有 1.20，范围 0--2，其中一个输入实现首次就没有脉冲。

这造成重要解释限制：一个原本接近零的读出，很难用相对下降推断“保留正常行为能力的选择性习惯化”。本研究没有建立脉冲计数到梳理发生概率、动作幅度或身体运动的映射，也没有直接输出刺激来检验执行能力。因此，主要结论应停留在 aBN1 的习惯化样神经响应，而不应写成“数字果蝇已经学会忽略触碰”。

\begin{table}[tbp]
\centering\small
\caption{下游读出暴露正常传递与地板问题。首次为第 1 次响应；末期先平均每例第 10--12 次，再跨五个输入实现平均。单位为脉冲数。}
\label{tab:downstream}
\begin{tabular}{lrrrr}
\toprule
模型 / 间隔 & aDN1 首次 & aDN1 末期 & aDN2 首次 & aDN2 末期\\
\midrule
\inputrows{data/downstream_rows.tex}
\bottomrule
\end{tabular}
\end{table}

\subsection{实际覆盖的习惯化特征}
\begin{table}[tbp]
\centering\small
\caption{习惯化特征与本轮证据。H 编号沿用经典综述的特征顺序\citep{rankin2009}；筛查阳性不等价于生物行为确证。}
\label{tab:hallmarks}
\begin{tabularx}{\linewidth}{p{3.7cm}Y}
\toprule
特征 & 当前证据状态\\
\midrule
H1 反应递减 & M1 短间隔 aBN1 达到预设门槛；M0 阴性；M1 长间隔效应较弱。\\
H2 自发恢复 & M1 短间隔通过神经输出筛查；使用独立恢复分支。\\
H3 跨轮次增强 & 未做多轮训练--恢复循环，未检验。\\
H4(a) 频率依赖递减 & 两间隔配对比较支持更频繁刺激导致更强递减。\\
H4(b) 频率依赖恢复 & 未建立；时间点稀疏、末态深度不同、平台条件不完全满足。\\
H5 强度依赖 & 仅测试一个输入率和脉冲宽度，未检验。\\
H6 渐近后的继续积累 & 未做专门的渐近后追加刺激对照。\\
H7 特异性/泛化 & 没有加入可比较的新刺激 B，未检验。\\
H8 去习惯化 & 没有 A--B--A 和相应时间/敏化对照，未检验。\\
H9 去习惯化本身递减 & 未检验。\\
H10 长期习惯化 & 未引入长期机制，也未进行长期协议。\\
\bottomrule
\end{tabularx}
\end{table}

\subsection{执行中断、复核和资源}
首次执行达到 1,800 s 工具时限后停止。其 manifest 保留 \texttt{interrupted} 状态，不把该批次伪装为完整成功。已完成的 M0 十条序列和 M1/1701 两条序列被保留；尚未完成的 M1/1702/T500 片段未进入汇总。随后两批按原参数补完剩余八条序列，合并覆盖检查无缺失、无重复。

有效序列所在检查点及批次的记录耗时合计约 3,065 s（51.1 min），包含加载、编译、快照、静默间隔计算和写盘，且不包含中断后未完成部分。因此它不是纯积分吞吐率，也不是所有实际占用时间的精确总和。两次接续记录的峰值 RSS 约 2.47 GiB；第一次中断未保存最终峰值，不能用接续值冒充所有批次的全局上限。

在整理本报告时，再次从原始事件运行独立分析器，重建的报告对象与原有结果完全一致。数据制表程序又将主表数值从逐种子指标重新计算，并逐项核对；没有重新模拟、拟合或挑选种子来改善论文图形。
\FloatBarrier

\section{讨论}
\label{sec:discussion}
\subsection{阴性结果说明了什么：正常传播与跨次记忆是不同能力}
M0 已经能够在 JO-CE 输入下产生正常回路响应，但在本轮两个间隔中没有逐次递减。这把“连线大致正确、能够传播信号”与“能在相应时间尺度保存刺激历史”区分开来。对初学者而言，关键不是模型有多少神经元，而是哪些状态能在两次刺激之间持续影响输出。

这一结果不能证明原模型在所有刺激、初态或参数下都没有慢动力学；也不能证明其他固定权重循环网络无法习惯化。我们没有扫描完整输入空间、系统搜索慢模态或改变原始参数。合理表述应始终附带条件：\emph{在所测试协议和参数范围内}，M0 不足以产生观察目标。

\subsection{充分性不是必要性，“最小扩展”不是唯一机制}
M1 表明，仅改变一组感觉输入的中央输出效能，就可以在固定图谱上得到所测的部分特征。这是一个模型内的\emph{充分性构造}：存在这样一组状态方程和参数，使现象出现。它不是\emph{必要性证明}，因为没有排除神经元适应、抑制增强、去抑制调节或其他慢状态机制。

$U=0$ 工程对照证明替换连接实现本身没有制造基线差异；它不等同于在真实动物中选择性阻断了某种分子机制，也不等于比较了所有候选记忆位置。M1 的作用点是 JO-CE 在脑内的兴奋性输出，而非经验证的习惯化特异性突触集合。只有进一步进行位置受约束的竞争模型与干预，才能回答机制在哪里、为什么在那里。

此外，M1 的强递减在一定程度上符合资源耗竭规则的直观预期。\textbf{“按会耗竭的规则建模，然后看到耗竭样下降”本身不足以构成新颖机制发现。}研究价值来自明确约束、阴性基线、实现对照，以及未来对未用于选择模型的协议和干预作出可被否定的预测。

\subsection{真实连接组的贡献尚未得到证明}
本轮保留了完整解剖结构，但没有比较裁剪子回路、随机图谱或低维模型。因此，不能把“在真实连接组上成功”改写成“只有真实连接组才能成功”。最小模型文献已经表明，小系统也可以展示下降与恢复\citep{smart2024}。

若要论证连接组有独特价值，应控制参数数量、拟合数据量和测试任务，并要求结构模型对留出刺激或细胞干预作出更好的预测。例如，两种模型都能拟合同一条 aBN1 曲线，但对抑制某个中间细胞的后果不同；若真实干预支持连接组模型，结构信息的贡献才得到更有力的证明。仅比较图的大小或视觉上更像大脑，不是有效对照。

\subsection{为什么恢复测试比单条下降曲线更有区分力}
短时耗竭、慢抑制和细胞适应都可能产生递减，但它们对休息长度、训练频率、追加刺激和旁路输入的预测可能不同。恢复测试因此有机会把“都能画出下降曲线”的模型分开。

本轮只有两个探测时间，且主读出是小整数计数，恢复曲线既稀疏又受量化影响。对于首次只产生少量脉冲的下游读出，少一个脉冲就可能导致很大的相对变化。这些因素使精确恢复速率估计不可靠。

还需避免把文献的数学结论直接移植到不同系统。Smart 等人的单状态最小模型具有明确的 H4(b) 限制\citep{smart2024}，但 M1 包含多个局部资源状态以及一个非线性全脑循环网络，并不是同一个数学对象。相同的局部恢复时间常数也不保证神经输出按同一个指数恢复。因此，本轮不把 H4(b) 自动判为通过或失败。

\subsection{正常功能损失与行为可解释性是实质性限制}
M1 的首次 aBN1 与下行响应降低，意味着“制造记忆”与“削弱整体传递”混在一起。一个好的生物候选模型不应仅在末期输出变小，还应在未训练时保持正常的刺激--输出关系，并在相关条件下保留应对新刺激的能力。

后续应预先设定正常功能门槛，例如要求多个下游读出在首次刺激时远离零计数地板；用直接激活下游或替代输入检验执行通路；同时观测输入细胞、中间层和输出层。若某个候选机制必须先把正常传递几乎关掉才能获得递减，就应把它视为重要失败模式，而不是只展示归一化后的漂亮曲线。

\subsection{本研究的主要局限}
\begin{enumerate}
\item \textbf{神经代理而非真实行为。}没有身体、运动反馈及脉冲到行为的校准。aBN1 计数不是梳理概率，aDN 发放也不等于动作执行。
\item \textbf{只测试一个输入强度和两种间隔。}没有覆盖弱/强刺激、更多时间尺度、其他模态及多个初态。
\item \textbf{同一图谱上的有限输入实现。}五个种子不能代表五个动物、基因型、个体连接差异或生物总体。
\item \textbf{冻结回放是一种控制而非生态真实性。}实际刺激往往逐次波动；本轮优先排除输入噪声混淆，尚未检验新的随机脉冲每次重抽时的稳健性。
\item \textbf{参数不是生理拟合结果。}$U$ 与恢复时间常数没有从 JO-CE 重复刺激数据估计，也未做敏感性或可辨识性分析。
\item \textbf{机制竞争与结构对照缺失。}没有证明 STD 优于抑制可塑性，也没有证明完整图谱优于小模型。
\item \textbf{关键习惯化特征缺失。}特异性、去习惯化、长期习惯化以及完整频率依赖恢复尚未确立。
\item \textbf{发表新颖性尚待审查。}本轮完成的是可重复计算先导研究，未进行穷尽性的新颖性检索，也没有足够证据作为完整生物机制论文直接投稿。
\end{enumerate}

\section{下一阶段：怎样把先导结果变成可发表的机制研究}
\label{sec:next}
以下为未来建议，不是本报告已经执行的实验。优先推进能区分假说的测试，而不是扩大参数搜索直到所有曲线都符合预期。

\subsection{竞争模型与判别实验}
建议至少比较三个候选：局部短时资源耗竭、中央抑制增强、抑制增强加解除抑制。各模型先通过相同的正常功能门槛，再用一部分协议约束少量共享参数；用从未参与拟合的协议评估预测。

\begin{table}[tbp]
\centering\small
\caption{建议的判别实验。预测须在新实验前由具体模型产生；下表不是已经测得的神经或行为结果。}
\label{tab:future}
\begin{tabularx}{\linewidth}{>{\raggedright\arraybackslash}p{3.2cm}YY}
\toprule
实验问题 & 设计要点 & 能排除或区分什么\\
\midrule
恢复有几种有效尺度？ & 从同一末态独立探测 0.2、0.5、1、2、5、10 s；检查平台与末态深度。 & 一种资源恢复尺度与多状态/多尺度模型的留出预测。\\
只是输出变弱了吗？ & 训练后直接驱动下游；检查未训练时的输出动态范围。 & 一般失能、低计数地板与保留潜在响应能力。\\
是否对刺激有选择？ & 选择能驱动可比较读出的 B，测试训练 A 对 A/B 的不同影响。 & 广泛疲劳与通路相关状态；不能使用无有效基线的 B。\\
B 是否解除 A 的递减？ & A--B--A、A--等时休息--A、未训练--B--A。 & 去习惯化、自然恢复与一般敏化。\\
抑制增强是否必要？ & 在明确标注的候选抑制细胞上设计选择性干预。 & 不同机制对旁路、沉默或增益改变的分歧。\\
真实结构是否有用？ & 同等数据预算比较最小模型、局部图、匹配随机图与全图。 & 结构带来的干预预测增益，而非仅有更高拟合自由度。\\
\bottomrule
\end{tabularx}
\end{table}

\subsection{什么时候应该转向嗅觉或视觉}
如果主问题转向刺激特异性和抑制性可塑性，嗅觉具有更直接的文献机制依据，但需补齐细胞标注、输入编码及分钟级时间尺度，不能把 JO-CE 参数直接照搬\citep{das2011}。如果主要目标是经典行为 hallmarks 与个体差异，light-off 跳跃更有成熟行为基础，但必须匹配视觉编码、相关输出及腹神经索/身体范围\citep{engel2009,boon2026}。

已有糖刺激基线可以成为跨模态工程对照，但唇瓣糖输入不能直接等同于足部糖的 PER 习惯化协议\citep{trisal2022}。选择模态应由要识别的机制和可获得的实验证据决定，而不是由最先跑通的演示锁定。

\section{结论}
在原始图谱、静态参数与本轮 JO-CE 协议下，固定连接全脑 LIF 没有出现跨次响应递减。仅对一组兴奋性输入输出边加入可恢复资源状态，即可在 aBN1 生成强递减、休息恢复和间隔依赖：短间隔的平均早晚期递减为 90.22\%，长间隔为 9.88\%。这些结果经过完整事件回读和输入配对校验。

最重要的结论不是“已经重建了会学习的果蝇”，而是：\textbf{正常感觉传递与跨次慢记忆必须分开检验；产生部分习惯化样现象的条件性充分性，并不等于真实机制识别或连接组必要性。}下一步应聚焦竞争机制、正常功能保留和留出干预预测。

\section*{数据、代码与研究透明性声明}
\addcontentsline{toc}{section}{数据、代码与研究透明性声明}
\textbf{代码与数据可得性。}本地项目包含固定协议、连续状态入口、独立分析器、逐种子结果、原始事件、输入实现及源码快照。已有公开仓库地址为：
\begin{center}\small\url{https://github.com/chenmzh/digital-fly-brain}\end{center}
本轮习惯化目录和本报告尚未推送；不能据该链接声称当前全部材料已公开。原始事件目前仅保留在本机，公开归档与长期数据托管仍待完成。

\textbf{伦理范围。}本轮只分析公开来源的处理后连接数据并进行计算仿真，没有新增活体动物或人体实验。原始数据的来源、使用条件与论文归属仍应遵守上游约定；这里不虚构伦理审批编号。

\textbf{署名、资助与利益冲突。}责任作者、单位、贡献分工、资助及利益冲突声明尚待人工确认，本文不据此虚构“无利益冲突”或任何资助信息。

\textbf{AI 辅助。}本项目的代码整理、文献归纳、分析辅助、制图和报告写作使用了 AI 编程助手。报告引用的数值来自实际运行和文件校验，但是否适合投稿、作者责任以及科学解释仍需人工审阅。

\FloatBarrier
\bibliographystyle{unsrtnat}
\bibliography{references}
\clearpage
\appendix
\section{术语与符号速查}
\label{app:glossary}
\small
\begin{longtable}{>{\raggedright\arraybackslash}p{3.1cm}>{\raggedright\arraybackslash}p{11.8cm}}
\caption{常用术语。类比只用于帮助阅读，不替代实际模型定义。}\label{tab:glossary}\\
\toprule
术语 & 本报告中的含义\\
\midrule\endfirsthead
\toprule 术语 & 本报告中的含义\\\midrule\endhead
\bottomrule\endfoot
神经元（neuron） & 能积累输入并在越过阈值时发出脉冲的模型单元；并未模拟真实细胞的全部生物过程。\\
脉冲（spike） & 一次离散神经发放事件。响应计数为指定窗口内的事件数，而非整段仿真的平均活动。\\
突触（synapse） & 神经元之间的作用接点。多处解剖接点可聚合为一条带权有向边。\\
连接组（connectome） & 细胞及其连接的结构图谱；不自动包含所有动力学参数或生前状态。\\
兴奋/抑制 & 本模型中增加或减少后续细胞发放倾向的传递符号；不等于“有用/有害”。\\
LIF & Leaky integrate-and-fire，泄漏积分发放。输入被积累，状态同时衰减，越阈值后发放并重置。\\
不应期（refractory period） & 发放后的短时特殊状态；本上游实现中暂停两个状态方程的更新。\\
习惯化（habituation） & 重复刺激后的反应递减及其相关性质。需与输入适应和输出失能区分。\\
去习惯化（dishabituation） & 新刺激之后，对原来已递减刺激的响应增加；不等于新刺激自身引起响应。\\
STD & Short-term synaptic depression，短时突触效能下降；不是“短时加强抑制性突触”的缩写。\\
JO-CE & 本实验驱动的 Johnston 器官机械感觉细胞群，固定清单含 70 个神经元。\\
aBN1 & 触角梳理相关脑内中间神经元，本研究主要读出。\\
aDN1/aDN2 & 相关下行神经元，本研究次要读出；不是完整身体模型。\\
读出（readout） & 从模型记录的观测量。本报告主要为细胞的窗口内脉冲数，而非动物行为标签。\\
冻结输入 & 同一种子生成后重复回放同一组事件，用来排除逐次输入随机变化。\\
状态快照 & 网络状态的保存副本，用于从同一训练末态产生独立恢复分支。\\
充分性/必要性 & “存在此机制就能产生现象”与“没有此机制就不能产生现象”是不同命题。\\
地板效应 & 初始计数已经很低甚至为零，进一步下降和比值就很难解释。\\
$T$ 与 $\Delta t$ & 前者是两次刺激起始的间隔；后者是仿真的数值时间步，二者相差多个数量级。\\
$I,E,L,D$ & 首次响应、早期均值、末期均值和递减比例。$D$ 的分母是 $E$，不是 $I$ 或 M0。\\
\end{longtable}
\normalsize

\section{逐种子结果与输入清单摘要}
\begin{table}[htbp]
\centering\small
\caption{M1 的全部逐种子主要结果。单位同主表；$D$ 为百分数。没有挑选“最好看”的种子。M0 在种子 1701--1704 的所有 aBN1 响应均为 10，在 1705 均为 9，因此不重复列出其恒定行。}
\label{tab:individual}
\begin{tabular}{rrrrrrrr}
\toprule
种子 & $T$(s) & $I$ & $E$ & $L$ & $D$(\%) & $P_2$ & $P_{10}$\\
\midrule
\inputrows{data/individual_rows.tex}
\bottomrule
\end{tabular}
\end{table}

\begin{table}[htbp]
\centering\small
\caption{每次刺激的注入事件与实际感觉神经元脉冲数。相同种子在全部对应条件下保持这些计数，且实际脉冲的相对时刻也一致。注入与发放的差额不是训练中逐次增大的效应；本轮未对每个差额事件的数值调度原因作单独归因。}
\label{tab:input}
\begin{tabular}{rrr}
\toprule
种子 & 注入事件数 & 实际感觉脉冲数\\
\midrule
\inputrows{data/input_rows.tex}
\bottomrule
\end{tabular}
\end{table}

\begin{table}[htbp]
\centering\small
\caption{读出细胞的 FlyWire v630 标识。JSON 清单采用字符串，原始事件 Parquet 使用 int64，避免浮点表示破坏大整数精度。}
\label{tab:ids}
\begin{tabular}{ll}
\toprule
读出 & 标识\\
\midrule
aBN1 & \texttt{720575940630907434}\\
aDN1 & \texttt{720575940616185531}\\
aDN2 & \texttt{720575940629806974}\\
\bottomrule
\end{tabular}
\end{table}
\FloatBarrier

\section{资源状态为何能够保存刺激历史：一个简短推导}
\label{app:derivation}
下面是对既定状态方程的解释，不是另一个已运行的实验，也不是生理参数拟合。没有事件到达时，式\eqref{eq:resource}有精确解
\begin{equation}
x(t)=1-\bigl[1-x(t_0)\bigr]\exp\left[-\frac{t-t_0}{\tau_{\mathrm{rec}}}\right].
\label{eq:recoveryanalytic}
\end{equation}
因此，在同一末态之后等待一个恢复时间常数，资源\emph{亏损}剩约 37\%；等待五个时间常数，亏损剩约 0.67\%。本研究人为设定 $\tau_{\mathrm{rec}}=2\secunit$，所以看到资源变量本身在 2 s/10 s 附近恢复，并不能算作独立发现。需要测量的是神经输出是否也恢复，以及这种关系能否预测新的刺激与干预。

为了理解频率效应，可把脉冲输入粗略近似为平均发放率 $f$，忽略具体到达时刻和网络反馈，得到
\begin{equation}
\frac{\mathrm{d}x}{\mathrm{d}t}\approx \frac{1-x}{\tau_{\mathrm{rec}}}-Ufx.
\end{equation}
若持续输入足够久，其近似平衡值为
\begin{equation}
x_\infty\approx\frac{1}{1+Uf\tau_{\mathrm{rec}}}.
\end{equation}
在 $U=0.01$、$f=150\,\mathrm{Hz}$、$\tau_{\mathrm{rec}}=2\secunit$ 时，该连续输入近似给出 $x_\infty=0.25$。\textbf{它不是本实验 200 ms 脉冲序列的实测平台值。}

有限脉冲之后，资源沿式\eqref{eq:recoveryanalytic}补充。若下一次刺激来得更早，就通常来不及补足；资源亏损可在多次刺激间积累。这解释了为什么单时间尺度机制可以产生间隔依赖。但从 $x$ 到 aBN1 的输出，还要经过异质连接、兴奋与抑制相互作用以及阈值发放，因此不能把输出计数简单当作 $x$ 的线性比例。

\section{复现、校验与审计路径}
\label{app:repro}
\subsection{主要文件及职责}
以下路径均相对于项目根目录。原始事件文件位于本地忽略目录，论文可读汇总与源文件位于可版本管理目录；两者不是同一种发布状态。
\begin{itemize}
\item \path{habituation/RESEARCH_PLAN.md}：正式结果前的研究问题、判据与停止条件。
\item \path{habituation/protocols/main_v1.json}：正式生物/数值协议。
\item \path{habituation/run_experiment.py}：连续状态全脑入口。
\item \path{habituation/analyze.py}：独立事件回读与指标计算。
\item \path{habituation/results/main_v1/}：中断批次的完整检查点及未纳入分析的片段。
\item \path{habituation/results/continuation_a/} 和 \path{habituation/results/continuation_b/}：同协议接续批次。
\item \path{habituation/analysis/responses.csv}：280 个响应窗口及辅助状态。
\item \path{habituation/analysis/metrics_by_seed.csv}：20 条序列的逐种子指标。
\item \path{habituation/analysis/report.json}：审计结果、来源哈希和描述性汇总。
\item \path{habituation/paper/prepare_data.py}：只读取已审计数据，生成本文制表与绘图输入。
\item \path{habituation/paper/report.tex}：本文 LaTeX 主文件；正文分节及文献在同目录。
\end{itemize}

\subsection{可执行复核流程}
从项目根目录，使用已有独立环境执行以下命令。新的输出目录必须不存在；不覆盖旧结果。
\begin{quote}\footnotesize\ttfamily
.venv/bin/python habituation/analyze.py\ \textbackslash\\
\hspace*{1em}--runs habituation/results/main\_v1\ \textbackslash\\
\hspace*{2em}habituation/results/continuation\_a\ \textbackslash\\
\hspace*{2em}habituation/results/continuation\_b\ \textbackslash\\
\hspace*{1em}--output habituation/results/new\_audit
\end{quote}
这会回读事件并要求完整的 20 条序列。只有主中断批次而没有接续批次时，分析器拒绝将其视为完整实验。工程单元测试覆盖恒定响应、递减后恢复、递减不恢复、零基线、平台、输入回放和非法协议，共 10 项；测试通过不等同于所有生物假设成立。

本轮正式协议的 SHA-256 为：
\begin{quote}\small\ttfamily
0c7477c0e5debaaf0a764ac6000285cb7f6\\
cbc8ef9e9fc468e7d889b3d8619f5
\end{quote}
两个印刷行直接连接构成完整的 64 位十六进制字符串。更多哈希保存在报告 JSON 和逐运行 manifest 中，而不是凭文字宣称“版本一致”。

\subsection{必须保留的偏离与工程事实}
\begin{enumerate}
\item \textbf{工具超时，不是阴性生物结果。}主运行被 1,800 s 工具时限终止，后续通过不改变参数的子集协议补完。中断片段保留但不参与统计。
\item \textbf{入口网格校验修复。}模拟启动后发现默认相对容差可能接受较长时间参数的半步错位，因此将校验改为绝对容差；正式协议本来就严格对齐网格，模型数值计算没有改变。历史入口快照与当前工作副本的区别被明确记录。
\item \textbf{探索性归一化。}“已损失响应的恢复比例”在看到部分结果后增加，只作为探索性诊断，未改写正式 H1/H2 判据。
\item \textbf{资源统计不夸大。}中断批次没有最终 RSS；记录耗时不是纯求解时间，也不是全部机器占用时间。
\item \textbf{没有追加模拟挑结果。}本报告排版阶段重新审计已有事件并生成图表，没有新增神经仿真、拟合参数或删掉不符合叙述的种子。
\end{enumerate}

\section{常见误读与阅读自检}
\begin{enumerate}
\item \textbf{“数字果蝇已经会学习了吗？”}本轮只能确认模型内部分神经响应特征，不能把它提升为完整动物学习或意识证据。
\item \textbf{“既然 M1 权重会变，为什么说不需要训练？”}基础解剖权重不经过优化器更新，但有效传递会随短时资源状态变化；短时可塑性与长期参数训练不是同一概念。
\item \textbf{“12 次刺激乘 5 个种子，就是 60 个独立样本吗？”}不是。同一序列内部状态相连；五个种子也只是同一图谱上的输入实现。
\item \textbf{“M1 通过 H1/H2，能叫动物习惯化吗？”}不能直接这样称呼。还缺少刺激特异性、输出能力与生物数据匹配等证据。
\item \textbf{“恢复均值等于首次，是否每例都精确恢复？”}不是。有些实现低于首次，有些高于首次，附录保留了这些差别。
\item \textbf{“M0 阴性意味着连接组没有用吗？”}不是。它在正常感觉传递中仍有作用；本轮只说明当前动态假设不足以产生所测试的慢历史效应。结构的独特贡献需要另一组对照来证明。
\item \textbf{“这是与 LLM 相同的学习吗？”}本研究使用真实连接约束的稀疏脉冲网络，显式推进毫秒级状态；没有语言 token、预训练目标或反向传播优化。两者都属于计算模型，但不能因都叫神经网络就等同其任务和证据标准。
\item \textbf{“论文体例就代表可以直接投稿了吗？”}不代表。署名、数据公开、机制竞争、预测验证与新颖性审查仍需完成。排版成熟度不能替代科学证据成熟度。
\end{enumerate}
\FloatBarrier

\end{document}
